\section{Evaluation and Results}
We evaluate each of the introduced components separately. Tables listing their interactions are available in the appendix. The appendix also lists some qualitative outputs.
The main trends that we observe are:
\begin{itemize}
    \item The new planner results in a small drop in BLEU, but is orders of magnitude faster (\S\ref{eval:planner}).
    \item Typing information causes a negligible drop in BLEU overall, but improves results substantially for the \emph{unseen} portion of the dataset (\S\ref{eval:typing}).
    \item The verification step is effective at improving the faithfulness of the output, practically eliminating omitted and overgenerated facts, reducing the number of wrong facts, and increasing the number of correctly expressed facts. This is based on both manual and automatic evaluations. (\S\ref{eval:verification}).
    \item The referring expression module is effective, with an internal correctness of \fix{@@@}. It substantially improves BLEU scores. (\S\ref{eval:reg}).
\end{itemize}
\begin{table}[t]

\begin{center}
\resizebox{\linewidth}{!}
{
\begin{tabular}{|l|l|l||l|l|}
\hline
 & \textbf{Moryosef et al} & \textbf{Moryosef et al} & \textbf{Exhaustive} & \textbf{Neural} \\
 & \textbf{StrongNeural} & \textbf{BestPlan} & \textbf{+Verify} & \textbf{+Verify} \\ \hline
\textbf{Expressed} & 360 & 417 & \textbf{426} & 405 \\ \hline
\textbf{Omitted} & 41 & 6 & \textbf{0} & 2 \\ \hline
\textbf{Wrong} & 39 & 17 & \textbf{14} & 30 \\ \hline
\textbf{Over-generation} & 29 & 3 & \textbf{0} & 4 \\ \hline
\textbf{Wrong REG} & - & - & \textbf{0} & 3 \\ \hline
\end{tabular}
}
\end{center}
\caption{Manual correctness analysis comparing our systems with the ones from \citet{step-by-step}.}
\label{table:manual}
\end{table}

\paragraph{Setup} We evaluate on the WebNLG dataset \cite{colin2016webnlg}, comparing to the step-by-step systems described in \citet{step-by-step}, which are state of the art. Due to randomness inherent in neural training, our reported automatic evaluation measures are based on an average of 5 training runs of each system (neural planner and neural realizer), each run with a different random seed. 

%We define a single experiment to be either using the different components presented in this paper, or using a parallel component from \cite{step-by-step} - Neural Planner vs Naive Planner, Token Typing vs No Typing, Output verification vs Best likelihood, and REG vs no REG.
%For each experiment, we perform 5 training runs with different random seeds for the trainable components - which are the neural planner and the neural realizer, and report the mean of each result.

\subsection{Neural Planner vs Exhaustive Planner}
\label{eval:planner}
\fix{@@table: Exhaustive planner vs. Neural Planner (the Base, Best K, No REG line.)}
\fix{@@table/figure: something about speed.}

We compare the exhaustive planner from \citet{step-by-step} to our neural planner, by replacing the planner component in the \citet{step-by-step} system. Table \ref{table:planning-bleu} shows that the move to the neural planner exhibits a small drop in BLEU. However, Table/@@Figure \ref{table:planning-speed} indicates an order of magnitude speedup.

\subsection{Effect of Typing Information}
\label{eval:typing}
@@Table: BLEU: exhaustive planner+types vs BLEU scores: all, seen, unseen
@@Table (maybe not needed): BLEU: neural planner+types vs BLEU scores: all, seen, unseen
The addition of typing information causes a small drop in BLEU overall, but substantially improve on the unseen part of the data. Additionally, for the unseen part of the data the verifier needed to explore 7.2 sentences on average before finding a verifiable one, compared to 5.8 with the typing information.

\subsection{Effect of Output Verification}
\label{eval:verification}
The addition of output verification resulted in negligible changes in BLEU, reinforcing that automatic metrics are not sensitive enough to output accuracy. We thus performed manual analysis, following the procedure in \citet{step-by-step}. We manually inspect @@ samples, containing 440 relations, counting expressed, omitted, wrong and over-generated (hallucinated) facts.\footnote{A wrong fact is one in which @@@, while an over-generated fact @@@.} We compare to the \textbf{StrongNeural} and \emph{BestPlan} systems from \citet{step-by-step}. Results in Table \ref{table:manual} indicate that the effectiveness of the verification process in ensuring correct output, reducing the already small number of ommited and overgenerated facts to 0 (when using the exhaustive planner) and keeping it small (when using the fast neural planner).     
@@put table here@@

\subsection{Referring Expression Module}
\label{eval:reg}
@@Table: the first line in current Table 1, split into two lines (one line for exhaustive planner, one for neural planner. The columns are - and REG).

\paragraph{Intrinsic evaluation of the REG module.}
We manually reviewed @@number@@ pairs of entities and referring expressions generated by the system.
We find that 92.2\% of the generated referring expressions refer to the correct entity.

Our of @@@ generated expressions, @@ (@@\%) were pronouns, @@ (@@\%) are repeating a one-token entity as is, and @@ (@@\%) are generating correct shortening of a long entity. In @@ (@@\%) of the cases the system did not find a good substitute and kept the entire entity intact. Finally, @@ (@@\%) are wrong referrals. Overall, 73.3\% of the non-first mentions of an entity were replaced with a suitable shorter and more fluent expressions.  
%We found the system manages to refer correctly 92.2\% of the time. Overall, 27.6\% are correct pronouns, 16.3\% are predicting to repeat the entity, given it contains one word, 2.25\% of the cases failed to find a valid replacement, therefore fell back on repeating the entire entity, and 7.8\% of the cases were wrong referrals. This concludes with 73.3\% of the repetitions being replaced with a more fluent expression.  

\paragraph{Effect on BLEU scores.}
As can be seen in Table \ref{@@@}, using the REG module increases BLEU scores for both the exhaustive and the neural planner.


\subsubsection{BLEU}
@@Move this to the appendix, with the title 'Interaction of different components' or something similar.

We analyze the output of our system using BLEU \cite{papineni2002bleu} on all 16 scenarios and report a per-scenario mean in Table \ref{table:bleu}.
In order to draw any conclusion if a specific component improves or worsen the BLEU score, we average all cells where that component appeared, and compare to the cells where that component didn't appear.
We find that using the verification component increases BLEU by 0.1, the REG component increases by 0.59, the Neural Planner worsen by 0.19, and the typing information worsen by 0.55.





\section{Conclusions}
We adopt the planning-based neural generation framework of \citet{step-by-step} and extend it to be orders of magnitude faster and produce more correct and more fluent text. We conclude that these extensions not only improve the system of \citet{step-by-step} but also highlight the flexibility and advantages of the step-by-step framework for text generation. 

\clearpage

\bibliography{emnlp-ijcnlp-2019}
\bibliographystyle{acl_natbib}
\clearpage
\appendix

\section{Neural Planner}
% \subsection{Further Details}
% We include the following pseudo-code to demonstrate how training is done.
% \algnewcommand{\IfThenElse}[3]{% \IfThenElse{<if>}{<then>}{<else>}
  \State \algorithmicif\ #1\ \algorithmicthen\ #2\ \algorithmicelse\ #3}
  
\begin{algorithm}
\caption{forward(g: Graph, plan: str[])}\label{code}
\begin{algorithmic}[1]
\State $\textit{edges} \gets \mathbf{E}(\mathbf{x_{n_i}};\mathbf{r_\ell};\mathbf{x_{n_j}})$
\State $\textit{nodes} \gets \mathbf{V}[\mathbf{x_{n_i}};\sum_{e_j\in \pi(i)}\mathbf{e_j};\sum_{e_j\in\pi^{-1}(i)}\mathbf{e_j}]$

\State $\textit{decoder} \gets LSTM().init(\overline{edges})$
\State $\textit{loss} \gets \emptyset$
\State $\textit{stack} \gets Stack()$

\While {$\textit{g.size} > 0$}
    \State $top \gets stack.top()$
    \If {$stack.empty$} 
        \State $\textit{options} \gets nodes$% \Comment{Options to traverse to a node}
    \Else
        \State $\textit{options} \gets \pi(t) \cup \pi^{-1}(t) \cup \{\vec{pop}\}$% \Comment{Options to traverse on an edge, or pop}
    \EndIf
    
    \State $y \gets plan[0]$
    \State $pred \gets options^{t} * decoder.out()$
    
    \State $loss.add(\mathcal{L}(pred, y))$
    \If {$\textit{node was selected}$}
        \State $decoder.add(node)$
        \State $stack.push(node)$
    \EndIf
    \If {$\textit{e=(s,r,o) was selected}$}
        \IfThenElse {$\textit{s} == \textit{top}$}{$\textit{d} \gets > $}{$\textit{d} \gets <$}
        \State $stack.push(o)$
        \State $decoder.add(d).add(r).add(``[").add(o)$
        \State $g.remove(e)$
    \EndIf
    \If {$\textit{pop was selected}$}
        \State $decoder.add(``]")$
        \State $stack.pop()$
    \EndIf
\EndWhile
\Return $\textit{loss}$
\end{algorithmic}
\end{algorithm}

\subsection{Speed}
We demonstrate the difference in speed between the naive-planner suggested in \citet{step-by-step}, and the neural-planner suggested in this paper in Figure \ref{fig:timing}. Figure \ref{fig:timing-naive} shows the naive-planner inference time is exponential and takes minutes to complete for the largest graphs in the corpus, and in contrast, Figure \ref{fig:timing-neural} shows the neural-planner inference time is linear and takes a few milliseconds to complete for the largest graphs in the corpus.

\begin{figure}[!ht]
    \begin{subfigure}{\linewidth}
        \includegraphics[width=\linewidth]{planning/timing/naive.png}
        \caption{Naive Planner}
        \label{fig:timing-naive}
    \end{subfigure}\\
    \begin{subfigure}{\linewidth}
        \includegraphics[width=\linewidth]{planning/timing/neural.png}
        \caption{Neural Planner}
        \label{fig:timing-neural}
    \end{subfigure}
    \caption{Error-line chart showing the time ($Y$ in seconds) on average it takes to create a plan for a graph with $X$ edge}
    \label{fig:timing}
\end{figure}

\section{Results}
For ablation, we ran every combination of components together to see the effect of each component. Our final system uses a combination of all of the components proposed in this paper.

% Table \ref{table:bleu} shows the BLEU score for an average of 5 independent runs for each combination.

\subsection{REG Error Analysis}
To elaborate on the types of errors, our unsupervised LM based REG component produces, we categorised all entities into 3 categories: (1) names of people; (2) locations (cities / counties / countries); and (3) places / objects.

We found that for the first category, our system produces no errors, and either selects a pronoun, the first name, or the last name of the person, all valid options.

For the second category, we found two distinct error categories. The first, when in the name of a location the name of the more specific location is not in BERT's vocabulary, for example ``Punjab, Pakistan'', it cannot predict the more specific word with only a single \emph{MASK} token, and thus selects ``Pakistan''. The second case is when the name of a city contains more than one word, for example ``New York''. While for ``New Jersey'' it is common to refer as ``Jersey'', it is unlikely to refer to ``New York'' with either ``New'' or ``York'', and as BERT can only fill in one MASK token, it chooses only one (in this case ``York'').

For the third category, we also found two distinctions. The first, when an entity contains more than one noun, it might be possible to select the correct one, like in ``Agra Airport'' in which we would like to say ``The Airport'', but not in other cases like ``Baked Alaska'', where choosing either word is not going to be a useful reference. 
The second case is specific for names of objects, like books, for example ``A Severed Wasp'', where we want the model to predict ``The Book''. However, as we don't allow words that are not pronouns, and are not in the original entity, this fails, and on this case, for example, predicts ``The Wasp''.

\subsection{Output Examples}
We add output examples to demonstrate the types of text our final system produces. The following outputs are correct, as they express all of the facts and only the facts from their input graphs. We enumerate them as number of facts:
\begin{enumerate}
    \item The leader of \textbf{Azerbaijan} is \textbf{Artur Rasizade}.
    \item \textbf{Baked Alaska}, containing \textbf{Sponge Cake}, is from \textbf{France}.
    \item \textbf{Above The Veil}, written by \textbf{Garth Nix}, is available in \textbf{Hardcover} and has \textbf{248} pages.
    \item The \textbf{Akita Museum Of Art} is located in \textbf{Japan} where the \textbf{Brazilians In Japan} are an ethnic group. \textbf{The Museum} is located in \textbf{Akita, Akita} which is part of \textbf{Akita Prefecture} .
    \item The \textbf{AWH Engineering College} in \textbf{Kuttikkattoor}, \textbf{Kerala} has \textbf{Mahé, India} to its northwest . \textbf{The College} was established in \textbf{2001} and has a staff of \textbf{250}.
\end{enumerate}

One example where the system failed, and wrongly lexicalized a fact is: ``The \textbf{AWH Engineering College} is located in the state of \textbf{Kerala}, \textbf{Kochi}, in \textbf{India}. The largest city in \textbf{India} is \textbf{Mumbai} and the river is the \textbf{Ganges}.''. In this example, the use of \textbf{Kochi} in the input is as the leader of \textbf{Kerala}, not as a location (although it is also one). The text lexicalizes this fact such that \textbf{Kerala} and \textbf{Kochi} are related, but with a relation of \emph{part-of}, implying \textbf{Kerala} is in \textbf{Kochi}.
